% Fragmento integrable. Preambulo anfitrion: amsmath, amssymb, longtable, hyperref.
\section{Identificación analítica de los centros originales}
\label{hmtfg:fragmento:identificacion-analitica-de-los-centros-originales}

\subsection{Realización de grado dos e identificación métrica de los centros}
\label{hmtfg:clasificacion:8}

La igualdad entre el primer parámetro del itinerario infinito y el cruce estable, demostrada por primera salida, proporciona
\begin{equation}
\boxed{a_*:=\min\Lambda_\tau=\lambda_*.}
\label{hmtfg:clasificacion:eq:H}
\end{equation}

La familia, el perfil dodecafásico, los retornos y el selector se mantienen exactamente como en secciones~\ref{hmtfg:clasificacion:1}--\ref{hmtfg:clasificacion:5}. Los resultados externos empleados a continuación reconocen la dinámica de ese objeto ya construido. Sus valores universales no intervienen en la selección de coeficientes ni de parámetros de la familia HMT.

\textbf{Teorema de transferencia métrica.} La identificación \eqref{hmtfg:clasificacion:eq:H}, junto con las cotas de entrada, admisibilidad y transversalidad anteriormente demostradas, implica que la sucesión original de primeras raíces nuevas satisface, para ciertas constantes \(C>0\) y \(0<\theta<1\),
\begin{equation}
\boxed{\lambda_*-\lambda_n=C\delta_{\mathrm F}^{-n}
\bigl(1+O(\theta^n)\bigr),}
\label{hmtfg:clasificacion:eq:14}
\end{equation}
donde \(\delta_{\mathrm F}>1\) es el autovalor inestable del operador de duplicación en el punto fijo reconocido. En particular,
\begin{equation}
\frac{\lambda_{n-1}-\lambda_{n-2}}
     {\lambda_n-\lambda_{n-1}}
=\delta_{\mathrm F}+O(\theta^n).
\label{hmtfg:clasificacion:eq:15}
\end{equation}
La sucesión coincide eventualmente con los centros de la continuación local, una vez que ambos índices se refieren al mismo período físico. La prueba mantiene el selector de primera raíz de sección~\ref{hmtfg:clasificacion:4}.

\subsubsection{Obtención de un dominio cuadrático semejante común}
\label{hmtfg:clasificacion:8.1}

Escribamos \(g\) para el punto fijo de la renormalización y
\(\Gamma(\lambda)=R^4f_\lambda\). La carta analítica utilizada en el certificado es
\[
f(x)=1-x^2V\!\left(\frac{x^2-1}{2.5}\right),\qquad
V(z)=\frac{u}{10}+\sum_{k\ge1}\nu_kz^k,
\]
con norma de diferencias \(|\Delta u|+\|\Delta\nu\|_1\). Las desigualdades de admisibilidad prueban que \(g\) es par, analítico, unimodal, cuadrático y de tipo de duplicación; además, \(g(x)=\varphi(x^2)\) con \(\varphi'<0\) en \([0,1]\).

El teorema 2.1 de de Faria–de Melo–Pinto, aplicado al conjunto de combinatorias formado únicamente por la duplicación, proporciona un conjunto límite de un solo punto con extensión cuadrática semejante. Como \(Rg=g\), la convergencia real de ese teorema identifica dicho punto con \(g\). El lema 3.2 proporciona, en un entorno analítico de \(g\), una iteración finita de renormalización con dominio cuadrático semejante y módulo uniformemente positivo. Los localizadores son secciones 2.1–2.3, pp. 736–738, y sección 3.1, lema 3.2, p. 745, de \href{https://annals.math.princeton.edu/wp-content/uploads/annals-v164-n3-p01.pdf}{de Faria–de Melo–Pinto, \emph{Global hyperbolicity of renormalization for \(C^r\) unimodal mappings}}.

Para comprobar la pertenencia al entorno de ese lema, sea \(\Omega_\epsilon\) un entorno complejo suficientemente pequeño de \([-1,1]\), con
\[
\rho_\epsilon:=\sup_{x\in\Omega_\epsilon}
\left|\frac{x^2-1}{2.5}\right|<1,
\qquad M_\epsilon:=\sup_{x\in\Omega_\epsilon}|x|^2<\infty.
\]
La restricción de la carta certificada satisface
\begin{equation}
\|\Delta f\|_{\Omega_\epsilon}
\le M_\epsilon\left(\frac{|\Delta u|}{10}
          +\rho_\epsilon\|\Delta\nu\|_1\right)
\le C_\epsilon\|\Delta f\|_{\mathcal A}.
\label{hmtfg:clasificacion:eq:16}
\end{equation}
El certificado da
\[
\|R^m\Gamma(\lambda_*)-g\|_{\mathcal A}
<0.001666(0.83996)^m.
\]
Se elige por ello un entero fijo \(m\) que sitúe esta imagen en el entorno del lema 3.2. La continuidad de la composición finita proporciona un intervalo \(I_*\) alrededor de \(\lambda_*\) donde las mismas operaciones están definidas. Si \(N\) es la iteración suministrada por dicho lema, ponemos
\begin{equation}
d=4+m+N,\qquad
\mathcal F_\lambda=R^df_\lambda,\quad \lambda\in I_*.
\label{hmtfg:clasificacion:eq:17}
\end{equation}
La familia \(\mathcal F\) posee representantes cuadráticos semejantes de grado dos, sobre un dominio común y con módulo positivo uniforme. Su dependencia analítica real admite una complejificación local del parámetro. El número \(d\) es fijo y cuenta las duplicaciones realizadas antes de aplicar el resultado paramétrico.

Esta construcción fija un dominio común y un número finito de retornos \(d\), suficientes para aplicar el teorema paramétrico.

\subsubsection{Transporte de la transversalidad certificada}
\label{hmtfg:clasificacion:8.2}

La dirección de \(\Gamma\) satisface el cono
\[
\|\dot\nu\|_1\le\tfrac14|\dot u|,\qquad \dot u\ne0.
\]
Las cotas diferenciales certificadas implican invariancia de este cono y
\begin{equation}
|\dot u_{k+1}|\ge4.381|\dot u_k|
\label{hmtfg:clasificacion:eq:18}
\end{equation}
a lo largo de la órbita de \(\Gamma(\lambda_*)\). Todos sus puntos permanecen en el dominio certificado. En particular, la dirección transportada por los pasos finitos de \eqref{hmtfg:clasificacion:eq:17} conserva una componente expansiva no nula.

La clase híbrida del punto fijo coincide con su variedad estable en el espacio de gérmenes cuadráticos semejantes: \href{https://www.maths.tcd.ie/EMIS/journals/Annals/149_2/lyubich.pdf}{Lyubich, \emph{Feigenbaum–Coullet–Tresser universality and Milnor's hairiness conjecture}, teorema 6.1, pp. 371–372}. La transversalidad sigue siendo válida al pasar a ese espacio, como puede comprobarse sin identificar normas de Banach diferentes. En efecto, si \(\partial_\lambda\mathcal F_{\lambda_*}\) fuese tangente a la clase híbrida, sería la tangente de un disco analítico contenido en ella. La convergencia exponencial uniforme de la renormalización sobre un subdisco compacto, seguida de la estimación de Cauchy en su parámetro, haría converger a cero la evaluación en \(x=1\) de sus tangentes renormalizadas. Sin embargo, en la carta anterior,
\[
\dot f(1)=-\dot u/10,
\]
y \eqref{hmtfg:clasificacion:eq:18} hace crecer su valor absoluto. Ambas derivadas representan el mismo germen y tienen la misma evaluación real. La contradicción demuestra
\begin{equation}
\partial_\lambda\mathcal F_{\lambda_*}
\notin T_{\mathcal F_{\lambda_*}}\mathcal H_{c_\infty}.
\label{hmtfg:clasificacion:eq:19}
\end{equation}
Así, \(S=\{\mathcal F_\lambda\}\) es una transversal analítica compleja a la clase híbrida del límite de duplicación.

\subsubsection{Identificación eventual de las primeras raíces por su período}
\label{hmtfg:clasificacion:8.3}

Denotemos por \(c_r\) el centro cuadrático canónico de período \(2^r\), obtenido mediante \(r\) duplicaciones sucesivas del centro de período uno. Este símbolo designa el parámetro de la familia cuadrática de reconocimiento; se distingue de las funciones orbitales \(c_j(\lambda)=f_\lambda^j(0)\) de las secciones anteriores. Se escribe \(c_r^{\rm quad}\) para conservar la distinción.

El teorema de continuación global y \eqref{hmtfg:clasificacion:eq:H} dan \(\lambda_n\to\lambda_*\), por lo que \(\lambda_n\in I_*\) para \(n\) suficientemente grande. La clasificación e intervalos restrictivos de sección~\ref{hmtfg:clasificacion:2} muestran que, cuando \(n>d\),
\begin{equation}
\mathcal F_{\lambda_n}=R^df_{\lambda_n}
\quad\text{tiene período crítico exacto }2^{n-d}
\quad\text{e itinerario }W_{n-d}0.
\label{hmtfg:clasificacion:eq:20}
\end{equation}
Su órbita crítica permanece en el intervalo real invariante, contenido en el dominio cuadrático semejante; su conjunto de Julia lleno es, por tanto, conexo. La rectificación conserva la órbita postcrítica y las ramas reales. Al recorrer sus retornos restrictivos se obtienen \(n-d\) duplicaciones canónicas y, al final, un punto crítico fijo. Este último rectifica al centro \(0\); las inversas sucesivas de la rectificación de las copias de duplicación dan el centro único \(c_{n-d}^{\rm quad}\). Esta es la aplicación concreta de la parametrización por combinatoria y del homeomorfismo de rectificación de las copias, descritos en Lyubich, secciones 5.1–5.3, pp. 362–364. Se concluye
\begin{equation}
\mathcal F_{\lambda_n}\in\mathcal H_{c_{n-d}^{\rm quad}}.
\label{hmtfg:clasificacion:eq:21}
\end{equation}

El teorema 7.4 de Lyubich, pp. 387–388, da una única intersección local de \(S\) con cada clase \(\mathcal H_{c_r^{\rm quad}}\) para \(r\) suficientemente grande. Sus hipótesis de cotas complejas a priori se cumplen para la cascada cuadrática real de duplicación, conforme al teorema 5.6, p. 367, del mismo trabajo. Escribamos \(\zeta_r\) para el parámetro de esa intersección. Las ecuaciones \eqref{hmtfg:clasificacion:eq:20}–\eqref{hmtfg:clasificacion:eq:21}, la convergencia al cruce y la unicidad local implican
\begin{equation}
\boxed{\lambda_n=\zeta_{n-d}\quad\text{para todo }n\text{ suficientemente grande}.}
\label{hmtfg:clasificacion:eq:22}
\end{equation}
La unicidad utilizada aquí corresponde a \textbf{cada clase híbrida canónica} en un entorno fijo del cruce. No se infiere de la mera existencia de una cascada local ni de la unicidad del cruce con la hoja estable. La ecuación \eqref{hmtfg:clasificacion:eq:22} identifica el selector original, ya definido globalmente, con las intersecciones locales; no introduce un selector sustitutorio.

En particular, sea \(\mu_j\) la cascada local de la sección \ref{hmtfg:escalado:8}, y sea \(s\) el entero fijo determinado por su período físico exacto:
\[
\operatorname{per}_{\rm crit}(f_{\mu_j})=2^{j+s}.
\]
En la presentación con curva \(R^4f_\lambda\) y blanco de período dos, \(s=5\); los pasos fijos adicionales usados para localizar el blanco modifican \(s\). Por \eqref{hmtfg:clasificacion:eq:21} y la misma unicidad,
\begin{equation}
\mu_j=\zeta_{j+s-d},\qquad
\boxed{\lambda_n=\mu_{n-s}\quad(n\text{ suficientemente grande}).}
\label{hmtfg:clasificacion:eq:23}
\end{equation}
El desplazamiento \(d\) de la preparación cuadrática semejante y el desplazamiento \(s\) de los períodos del blanco son objetos distintos. La identificación se efectúa por el período físico, no igualando esos dos enteros.

\subsubsection{Resto de potencia mediante holonomía transversal}
\label{hmtfg:clasificacion:8.4}

El lema 7.3 de Lyubich, p. 384 y demostración en pp. 384–387, aporta regularidad \(C^{1+\beta}\)-conforme, para algún \(\beta>0\), de la holonomía entre transversales a la clase híbrida del punto de Feigenbaum. Su definición en p. 384 se refiere a los conjuntos de conexidad sobre las transversales; contiene, en particular, los centros considerados aquí. Podemos reducir el exponente para que \(0<\beta\le1\).

Sea \(h\) la holonomía desde \(S\) hacia la variedad inestable y sea \(z\) una coordenada analítica que linealiza la restricción unidimensional de la renormalización:
\begin{equation}
z(g)=0,\qquad z(Rq)=\delta_{\mathrm F} z(q).
\label{hmtfg:clasificacion:eq:24}
\end{equation}
La existencia de esta coordenada se obtiene aplicando la linealización de un germen holomorfo repulsor, de multiplicador \(\delta_{\mathrm F}>1\), a la variedad inestable analítica. La dirección expansiva de \eqref{hmtfg:clasificacion:eq:18}, la invariancia del cono y el retorno estacionario identifican ese multiplicador con el autovalor inestable positivo empleado en el certificado.

Si \(q_r\) es el centro de la clase \(\mathcal H_{c_r^{\rm quad}}\) situado en la variedad inestable, la compatibilidad de rectificación y renormalización da \(Rq_r=q_{r-1}\). Esta identidad es la que se utiliza en la demostración del teorema 7.4, p. 388. Por \eqref{hmtfg:clasificacion:eq:24}, existe \(b\ne0\), que absorbe cualquier índice inicial fijo, tal que
\begin{equation}
z(q_r)=b\delta_{\mathrm F}^{-r}
\label{hmtfg:clasificacion:eq:25}
\end{equation}
para todo \(r\) suficientemente grande.

La regularidad de \(h\) y la transversalidad \eqref{hmtfg:clasificacion:eq:19} dan, sobre el conjunto de conexidad cercano a \(\lambda_*\),
\begin{equation}
H(\lambda):=z\bigl(h(\mathcal F_\lambda)\bigr)
=A(\lambda-\lambda_*)
+O\!\left(|\lambda-\lambda_*|^{1+\beta}\right),
\qquad A\ne0.
\label{hmtfg:clasificacion:eq:26}
\end{equation}
Como \(H(\zeta_r)=b\delta_{\mathrm F}^{-r}\), la estimación inferior
\(|H(\lambda)|\ge |A|\,|\lambda-\lambda_*|/2\) en un entorno suficientemente pequeño implica primero
\(|\zeta_r-\lambda_*|=O(\delta_{\mathrm F}^{-r})\). Sustituyendo esta cota en el resto de \eqref{hmtfg:clasificacion:eq:26}, resulta
\begin{equation}
\zeta_r-\lambda_*=\frac bA\delta_{\mathrm F}^{-r}
+O\!\left(\delta_{\mathrm F}^{-(1+\beta)r}\right).
\label{hmtfg:clasificacion:eq:27}
\end{equation}
Aplicamos \eqref{hmtfg:clasificacion:eq:22} con \(r=n-d\). El desplazamiento finito modifica el coeficiente principal y la constante del resto. Puesto que \(\lambda_n<\lambda_*\), se obtiene \eqref{hmtfg:clasificacion:eq:14}, con
\[
C=-\frac bA\delta_{\mathrm F}^d>0,
\qquad \theta=\delta_{\mathrm F}^{-\beta}\in(0,1).
\]
Finalmente,
\[
\lambda_n-\lambda_{n-1}
=C(\delta_{\mathrm F}-1)\delta_{\mathrm F}^{-n}\bigl(1+O(\theta^n)\bigr),
\]
lo que demuestra \eqref{hmtfg:clasificacion:eq:15}. \(\square\)


\subsubsection{Alcance cuantitativo del teorema}
\label{hmtfg:clasificacion:8.5}

La identificación de límites demostrada por primera salida activa la transferencia métrica para las primeras raíces globales del sector uniforme. La tasa operatorial \(0.83996\) controla los mapas renormalizados sobre la hoja estable. La tasa paramétrica \(\theta=\delta_{\mathrm F}^{-\beta}\) procede de la regularidad de la holonomía transversal. Ambas estimaciones conservan su variable y su dominio.

El teorema acredita la existencia de \(\beta,\theta,C\), de la vecindad \(I_*\), del número finito de retornos \(d\) y de un índice \(n_0\) a partir del cual coinciden los centros. Sus valores numéricos individuales no se asignan en este enunciado asintótico. Esta formulación completa el límite y su resto de potencia; la evaluación numérica de esos testigos constituye una especialización cuantitativa distinta.
